% Edited conversion of source p. 30.
\ReportHeading
  {Дослідження пружнопластичного стану тришарових циліндричних оболонкових елементів ракетно-космічної техніки з отворами методом скінченних елементів}
  {Є.~А. Сторожук, Р.~В. Серафимович, І.~С. Чернишенко}
  {Інститут механіки ім. С.~П. Тимошенка НАН України}
  {}

Шаруваті пластини та оболонки широко застосовують в авіа- і суднобудуванні, космічній техніці, радіоелектроніці, промисловому та цивільному будівництві. Порівняно з одношаровими елементами такої самої маси шаруваті конструкції можуть мати вищі жорсткість і міцність. Відповідний вибір матеріалів зовнішніх шарів та заповнювача забезпечує додаткові властивості: звуко- і теплоізоляцію, вібропоглинання та теплостійкість.

Значного поширення набули тришарові циліндричні оболонки. Через конструктивні або технологічні вимоги вони часто мають отвори й вирізи різної форми. До таких елементів належать перехідні, приладові, хвостові та паливні відсіки ракет-носіїв. За підвищених навантажень біля отворів виникають зони концентрації напружень, а властивості матеріалів шарів можуть описуватися нелінійними діаграмами деформування.

Більшість відомих результатів щодо розподілу напружень навколо отворів у тришарових циліндричних оболонках отримано для лінійно-пружної стадії. Систематичні дослідження фізично нелінійного деформування таких оболонок залишаються обмеженими.

Запропоновано чисельну методику розрахунку концентрації напружень біля отворів у тришарових циліндричних елементах ракетно-космічної техніки з урахуванням нелінійних властивостей матеріалів шарів. Методика ґрунтується на гіпотезі ламаної лінії, теорії малих пружнопластичних деформацій, методі додаткових напружень і методі скінченних елементів~\cite{sandwichholes:kudin2019,sandwichholes:reddy2004,sandwichholes:guz1980,sandwichholes:fem1982}. Для заповнювача використано співвідношення теорії Тимошенка, а для несучих шарів --- теорії Кірхгофа--Лява.

За допомогою розробленої методики та програмного забезпечення досліджено пружнопластичний стан тришарового циліндричного перехідного відсіку з круговим отвором під дією осьового навантаження.

\begin{thebibliography}{9}
\bibitem{sandwichholes:kudin2019}
О.~В. Кудін, Є.~А. Сторожук, С.~В. Чопоров,
\emph{Наближені аналітичні та чисельні методи аналізу міцності тришарових тонкостінних конструкцій}.
Херсон: Гельветика, 2019, 160~с.

\bibitem{sandwichholes:reddy2004}
J.~N. Reddy,
\emph{Mechanics of Laminated Composite Plates and Shells: Theory and Analysis}.
Boca Raton: CRC Press, 2004, 568~p.

\bibitem{sandwichholes:guz1980}
А.~Н. Гузь, И.~С. Чернышенко, Вал.~Н. Чехов, Вал.~Н. Чехов, К.~И. Шнеренко,
\emph{Методы расчета оболочек. Т.~1: Теория тонких оболочек, ослабленных отверстиями}.
Киев: Наукова думка, 1980, 636~с.

\bibitem{sandwichholes:fem1982}
\emph{Метод конечных элементов в механике твердых тел}, под ред. А.~С. Сахарова и И.~Альтенбаха.
Киев: Вища школа, 1982, 480~с.
\end{thebibliography}
